\begin{solapartat}
\punts{0,4} $f = \dfrac{c}{\lambda}$

\punts{0,5} $E_{\textit{fotó}} = hf$

\punts{0,35}

\begin{center}
\begin{tabular}{cccc}
\toprule
$\lambda$ ($\mu$m) & $f$ (Hz) & $E$ (J) & $E$ (eV)\\
\midrule
1,04 & $2,88\times 10^{14}$ & $1,91\times 10^{-19}$ & 1,19\\
0,6 & $5\times 10^{14}$ & $3,32\times 10^{-19}$ & 2,07\\
0,5 & $6\times 10^{14}$ & $3,98\times 10^{-19}$ & 2,48\\
\bottomrule
\end{tabular}
\end{center}
\end{solapartat}

\begin{solapartat}
\punts{0,5} $E_{C,\textit{màx}} = hf - W_e$

\begin{center}
\begin{tabular}{ccc}
\toprule
$\lambda$ ($\mu$m) & $E$ (eV) & $E_{C,\textit{màx}}$ (eV)\\
\midrule
1,04 & 1,19 & -\\
0,6 & 2,07 & 0,87\\
0,5 & 2,48 & 1,28\\
\bottomrule
\end{tabular}
\end{center}

\punts{0,25} La freqüència llindar correspon a la dels fotons que tenen una energia $W_e$:
\[ f = \frac{W_e}{h} = \frac{1,2\cdot 1,602\times 10^{-19}}{6,63\times 10^{-34}} = 2,90\times 10^{14}\ \mathrm{Hz} \]
No hi ha electrons arrencats per $\lambda_1$ atès que la seva freqüència, $2,88\times 10^{14}$~Hz, està per sota del llindar.

També és pot argumentar que no s'arrenquen electrons per $\lambda_1$ atès que la seva energia, 1,19~eV, és inferior a $W_e$.

\figura[0.5]{sol_fig1.png}

\punts{0,25} Escalat eixos correcte, punts correctament representats.

\punts{0,25} Eixos amb títols i unitats.

Si es representa l'extrapolació de les dades de manera que s'entra en la regió d'energies negatives, llavors es restarà 0,2~p a la qualificació.
\end{solapartat}
